\documentstyle[% 


righttag% 


,11pt% 


,amssymb% 


,russian%               remove in Eng. 


,izvru%                 Set izven% in Eng. 


% ,izvru-ks             for brief communications 


]{amsart} 


 


% 


 


\newcommand{\mes}{\operatorname{mes}} 


\newcommand{\Div}{\operatorname{div}} 


\newcommand{\tts}[1]{} 


 


%\hoffset=-15mm%                  For Eng. 


\setcounter{page}{58} 


\god{2004} 


\nomer{12~(511)} %                        Remove in Eng. 


%\nomer{Vol.\,48, No.\,12}%               Set in Eng. 


%\str{\pageref{begin}--\pageref{end}}%  Set in Eng. 


\udk{517.977} 


 


\author{\it С.Я.\,Серовайский} 


% NB:   ^^ remove in Eng. 


\title{Оптимальное управление для сингулярного уравнения 


с негладким оператором и изопериметрическим условием} 


 


\date{} 


%\pagestyle{myheadings}%         Set in Eng. 


\pagestyle{plain} %              Remove in Eng. 


\begin{document} 


%\thispagestyle{plain}%          Set in Eng. 


\maketitle 


\label{begin} 


 


 


 


 


\section{Постановка задачи} 


 


Задана открытая ограниченная область $\Omega$ пространства $\Bbb R^n$. В 


дальнейшем для краткости через $L_p$, $H^1_0$ и т.\,д. будут 


обозначены пространства $L_p(\Omega)$, $H^1_0(\Omega)$ и т.\,д. В 


соответствии с 


теоремой Соболева имеют место непрерывные вложения $H^1_0\subset L_q$, 


$L_{q'}\subset H^{-1}$, где 


$1/2-1/n=1/q$, $n>2$, $q$ произвольно при $n=2$, $1/q+1/q'=1$. 


Рассматривается уравнение 


\begin{equation} 


\Delta y+g(y)=v+f 


\end{equation} 


с однородным граничным условием (задача Дирихле), где $v$ --- 


управление, $f$ --- известная функция из пространства $H^{-1}$, $g$


--- заданная функция от состояния системы $y$. Управление $v$ 


выбирается из выпуклого замкнутого подмножества $U$ пространства 


$L_2$. Предполагается, что функция $g$ принадлежит классу $G$ 


непрерывных функций, удовлетворяющих неравенству $|g(y)|\le c|y|^{q-1}$ 


$\forall y$, где 


через $c$ здесь и далее будут обозначаться различные 


положительные константы. Тогда, пользуясь теоремой 


Красносельского ([1], с.\,312), установим, что оператор 


Немыцкого $g:L_q\to L_{q'}$ является непрерывным. Если бы на функцию $g$ 


были 


наложены стандартные ограничения монотонности и коэрцитивности, 


то с помощью теории монотонных операторов (напр., [2], 


с.\,184) можно было показать, что для любого управления $v\in U$ 


задача Дирихле для уравнения (1) имеет единственное решение $y$ 


из $H^1_0$. Однако в данном случае мы не имеем возможности 


воспользоваться этим результатом. Более того, по крайней мере, 


для некоторых значений функции $g$ однозначная разрешимость для 


уравнения (1) не имеет места (напр., [3], с.\,262), что, 


впрочем, не является препятствием для постановки и 


исследования оптимизационных задач. 


 


В соответствии с общей концепцией решения экстремальных задач 


для сингулярных уравнений (напр., [3], [4]) под множеством $W$ 


допустимых пар для уравнения (1) будем понимать пару функций $(v,y)$ 


из множества $S_U=U\times H^1_0$, удовлетворяющих этому уравнению. 


Задается функционал 


$$ 


K(y)=\int_\Omega k[x;y(x)]dx, 


$$ 


где $k$ --- функция Каратеодори. Пусть множество $W_\partial=\{(v,y)\in 


W\mid K(y)=0\}$ не пусто. 


Рассматривается функционал 


$$ 


I(v,y)=\int_\Omega F[x;y(x),\nabla y(x)]dx+ 


\frac{\gamma}{2}\|v\|^2_2, 


$$ 


где $F$ --- функция Каратеодори на $\Omega\times\Bbb R^{n+1}$, 


удовлетворяющая 


неравенству $F(x;y,z)\ge c|z|^2$ почти всюду (п.\,в.) на $\Omega$ $\forall 


y\in\Bbb R$, $z\in\Bbb R^n$, $\gamma>0$, через $\|\cdot\|_p$ 


обозначается норма в пространстве~$L_p$. Ставится оптимизационная 


\medskip 


 


{\it Задача $P$}. Найти такую допустимую пару $(v,y)$, которая 


минимизирует функционал $I$ на множестве $W_\partial$. 


 


\begin{thm} 


Если справедливо неравенство $|k(x,y)|\le a(x)+cy^2$, где $a\in L_1$, а 


функция $F(x;y,\cdot)$ выпукла 


почти для всех $x\in\Omega$ и $y\in\Bbb R$, то задача $P$ разрешима. 


\end{thm} 


 


\begin{pf} 


Поскольку минимизируемый функционал ограничен снизу, он 


обладает нижней гранью на множестве $W_\partial$. Пусть 


последовательность $\{w_k\}=\{v_k,y_k\}$ допустимых пар является 


минимизирующей, 


т.\,е. справедливы включение $w_k\in W_\partial$, равенство 


\begin{equation} 


\Delta y_k+g(y_k)=v_k+f 


\end{equation} 


и сходимость $I(w_k)\to\inf I(W)$. Из определения минимизируемого 


функционала 


следует ограниченность последовательности $\{w_k\}$ в пространстве 


$S=L_2\times H^1_0$. Тогда после выделения подпоследовательностей (с 


сохранением прежнего обозначения) имеем сходимость $w_k\to w=(v,y)$ слабо 


в $S$. Учитывая выпуклость множества $U$, установим включение $v\in U$. 


В силу ограниченности последовательности $\{y_k\}$ в пространстве $L_q$ 


и определения множества $G$ последовательность $\{g(y_k)\}$ ограничена в 


пространстве $L_{q'}$. Тогда в результате извлечения 


подпоследовательности установим, что $g(y_k)\to g$ в $L_{q'}$. Отметим, 


что последовательность $\{f_k\}$, где $f_k=v_k+f-g(y_k)$, оказывается 


ограниченной в 


пространстве $L_{q'}$. Отсюда следует, что решение $y_k$ однородной 


задачи Дирихле для уравнения Пуассона $\Delta y_k=f_k$ ограничено в 


пространстве $W^2_{q'}$. Применяя теорему Реллиха--Кондрашова, после 


выделения подпоследовательности получаем сходимость $y_k\to y$ сильно 


в $L_2$ и п.\,в. на $\Omega$, а значит, $g(y_k)\to g(y)$ п.\,в. на 


$\Omega$. Пользуясь 


леммой 1.3 ([2], с.\,25), установим, что $g=g(y)$. При заданных 


ограничениях на функцию $k$ из теоремы Красносельского следует 


непрерывность функционала $K$ в пространстве $L_2$. Переходя к 


пределу в равенстве $K(y_k)=0$, получаем $K(y)=0$. После перехода к 


пределу 


в равенстве (2) заключаем, что $w\in W$. Пользуясь теоремой о 


полунепрерывности снизу интегрального функционала ([5], 


теорема 2.1, с.\,242), получим неравенство $I(w)\le\varliminf I(w_k)$. 


Следовательно, 


пара $w$ является решением задачи $P$. 


\end{pf} 


 


 


Вопросы оптимального управления системами, описываемыми 


нелинейными уравнениями эллиптического типа с однозначно 


разрешимыми уравнениями состояния, рассматриваются, например, 


в [6]--[8]. Распространение этих результатов на задачу $P$ не 


представляется возможным в виду того, что в процессе 


варьирования управления не гарантировано получение допустимой 


пары. В работах [3], [4], посвященных решению сингулярных 


бесконечномерных оптимизационных задач, рассматриваются 


гладкие операторы. Описанные там подходы можно было бы 


использовать в сочетании с известными методами негладкой 


оптимизации (напр., [5], [9]), если бы негладкие члены входили 


в критерий оптимальности, а не в уравнение состояния. Вопросы 


оптимального управления системами с фазовыми ограничениями 


рассматриваются, например, в [10]--[12]. Однако эти 


исследования не адаптированы к сингулярным бесконечномерным 


системам. Желаемый результат в данном случае получается за 


счет модификации введенного в [13] понятия слабого 


приближенного решения оптимизационной задачи и описанной в 


[14] гладкой аппроксимации уравнения. Отметим, что 


исследование в дальнейшем будет проводиться без ограничений на 


выпуклость подинтегральной функции в критерии оптимальности и 


на степень роста функции $k$, т.\,е. при нарушении условий 


теоремы 1. Тем самым задача $P$ может оказаться неразрешимой, 


что также препятствует применению методов, описанных в работах 


[6]--[12]. В этих условиях получаемые результаты укладываются 


в общую концепцию расширения экстремальных задач (напр., [5], 


[15], [16]), хотя используемый принцип расширения здесь 


качественно иной. 


 


Определим понятие приближенного решения задачи $P$. 


Последовательность $\{w_n\}=\{v_n,y_n\}$ пространства $S$ назовем {\it 


слабой 


минимизирующей последовательностью} для задачи $P$, если для 


любых $\varepsilon>0$, $\delta>0$ и любой окрестности $O$ нуля в слабой 


топологии 


этого пространства существует такая допустимая пара $w'$ системы 


(1) и такой номер $n(\varepsilon,\delta,O)$, что справедливы соотношения 


$v_n\in U$,


$w_n\in w'+O$, $|K(y_n)|\le\delta$ и $I(w_n)\le\inf I(W)+\varepsilon$ 


$\forall n>n(\varepsilon,\delta,O)$. Если известна слабая минимизирующая 


последовательность, 


то, выбирая номер $n$ достаточно большим, можно обеспечить 


сколь угодно высокую степень близости величины $w_n$ к множеству 


допустимых пар системы (1) (эта точка попадает в сколь угодно 


малую окрестность некоторой допустимой пары), выполнение 


изопериметрического условия в точке $y_n$ с произвольной степенью 


точности, а также близость значения $I(w_n)$ к нижней грани 


минимизируемого функционала на этом множестве. Соответствующую 


пару $w_n$ будем называть {\it слабым приближенным решением}


задачи~$P$. Дадим способ его построения с помощью некоторого метода 


аппроксимации. 


 


 


\section{Аппроксимация задачи} 


 


Для решения поставленной задачи осуществляется гладкая аппроксимация 


уравнения состояния. Рассмотрим последовательность функций $\{g_n\}$ 


класса $G$, 


удовлетворяющих условиям 


\begin{equation} 


g_n\in C^1(\Bbb R),\ \ n=1,2,\dotsc;\ \ g_n(y)\to g(y) 


\ \text{равномерно по}\ y\in\Bbb R. 


\end{equation} 


Аналогично определяется последовательность $\{k_n\}$ непрерывно 


дифференцируемых по второму аргументу функций, удовлетворяющих условиям 


\begin{equation} 


|k_n(x,y)|\le k_0(x),\ \ n=1,2,\dotsc;\ \ k_n(x,y)\to k(x,y) 


\ \text{равномерно по}\ y\in\Bbb R 


\end{equation} 


почти для всех $x\in\Omega$, где функция $k_0$ интегрируема. Определим 


функционал 


$$ 


I_n(v,y)=I(v,y)+\frac{1}{q'\varepsilon_n}\|\Delta y+g_n(y) 


-v-f\|^{q'}_{q'}+\frac{1}{2\varepsilon_n}\bigg|\int_\Omega 


k_n[x,y(x)]dx\bigg|^2+\frac{\varepsilon_n}{2} 


\|\Delta y\|^2_2, 


$$ 


где $\varepsilon_n>0$ и $\varepsilon_n\to0$ при $n\to\infty$. 


 


Рассматривается экстремальная 


\medskip 


 


{\it Задача $P_n$}. Найти такую пару, которая минимизирует функционал 


$I_n$ на множестве $S_U$. 


\medskip 


 


Отметим, что второе слагаемое в функционале $I_n$ вводится в 


соответствии с общей концепцией метода штрафа [3]. Однако в 


отличие от его стандартной версии здесь осуществляется также 


аппроксимация, вообще говоря, негладких членов уравнения 


состояния. Аналогичный смысл имеет и третье слагаемое, 


реализующее ``аппроксимационный штраф'' изопериметрического 


условия. Кроме того, с целью получения дополнительной 


априорной оценки на последующей стадии исследования функционал 


дополняется четвертым слагаемым. В определенной степени оно 


напоминает стабилизатор из метода регуляризации Тихонова, хотя 


обычно при этом используют норму управления [17]. По аналогии 


с теоремой 1 доказывается 


 


\begin{thm} 


Для любого номера $n$ задача $P_n$ разрешима. 


\end{thm} 


 


\begin{pf} 


Пусть последовательность $\{w_k\}=\{v_k,y_k\}$ допустимых пар является 


минимизирующей, т.\,е. справедливы включение $w_k\in S_U$ и сходимость 


$I_n(w_k)\to\inf I_n(S_U)$. Из определения функционала $I_n$ следует 


ограниченность 


последовательности $\{w_k,\Delta y_k,h_k,m_k\}$ в пространстве $S\times 


L_2\times L_{q'}\times\Bbb R$, где 


\begin{equation} 


\Delta y_k+g_n(y_k)-v_k-f=h_k,\quad m_k=\int_\Omega 


k_n[x,y_k(x)]dx. 


\end{equation} 


Из определения множества $G$ следует ограниченность $\{g_n(y_k)\}$ в 


$L_{q'}$. 


Тогда после выделения подпоследовательностей имеем сходимость 


$w_k\to w=(v,y)$ слабо в $S$, $\Delta y_k\to\Delta y$ слабо в $L_2$, 


$h_k\to h$ слабо в $L_{q'}$, $g_n(y_k)\to g$ слабо в $L_{q'}$, 


$m_k\to m$, причем $v\in U$. Извлекая подпоследовательность, установим, 


что 


$y_k\to y$ сильно в $L_2$ и п.\,в. на $\Omega$, а значит, $g_n(y_k)\to 


g_n(y)$, $k_n(x,y_k)\to k_n(x,y)$, а также $\nabla y_k\to \nabla y$ п.\,в. 


на $\Omega$. Отсюда следует $g=g_n(y)$. В результате перехода к пределу 


в первом равенстве (5) заключаем, что $\Delta y+g_n(y)-v-f=h$. Применяя 


теорему 


Лебега о мажорированной сходимости ко второму равенству (5), 


установим 


$$ 


m=\int_\Omega k_n[x,y(x)]dx. 


$$ 


Тогда, пользуясь теоремой 2.1 ([5], с.\,242), получаем соотношение 


$I_n(w)\le\varliminf I_n(w_k)$, 


а значит, пара $(v,y)$ действительно является решением задачи $P_n$. 


\end{pf} 


 


Для обоснования сходимости метода регуляризации 


предполагается, что скорость сходимости последовательности 


$\{\varepsilon_n\}$ 


ниже, чем у $\{g_n\}$ и $\{k_n\}$ в том смысле, что выполняются 


соотношения 


\begin{equation} 


\lim_{n\to\infty}\frac{1}{\varepsilon_n}\sup_{y\in\Bbb R} 


|g_n(y)-g(y)|^{q'}=0,\quad \lim_{n\to\infty} 


\frac{1}{\sqrt{\varepsilon_n}}\int_\Omega \sup_{y\in\Bbb R} 


|k_n(x,y)-k(x,y)|dx=0. 


\end{equation} 


 


\begin{thm} 


Последовательность $\{w_n\}=\{v_n,y_n\}$ решений задачи $P_n$ является 


слабой минимизирующей последовательностью для задачи $P$. 


\end{thm} 


 


\begin{pf} 


В силу ограниченности снизу функционала $I$ на множестве $W_\partial$ 


существует его нижняя грань на этом множестве. Тогда для 


произвольного числа $\delta>0$ найдется такая точка $w=(v,y)$ из множества 


$W_\partial$, чтобы выполнялись соотношения $I(w)\le\inf 


I(W_\partial)+\delta$. В результате установим 


соотношение 


\begin{gather*} 


I_n(w_n)=\min I_n(S_U)\le I_n(w)=I(w)+\frac{\varepsilon_n}{2} 


\|\Delta y\|^2_2+ \\ 


% 


+\frac{1}{q'\varepsilon_n}\int_\Omega|g_n(y(x))-g(y(x))|^{q'} 


dx+\frac{1}{2\varepsilon_n}\bigg|\int_\Omega 


\{k_n[x,y(x)]-k[x,y(x)]\}dx\bigg|^2. 


\end{gather*} 


Таким образом, справедливо неравенство 


\begin{gather*} 


I_n(w_n)\le I(w)+\frac{\varepsilon_n}{2}\|\Delta y\|^2_2+\\ 


% 


+\frac{1}{q'\varepsilon_n}\sup_{y\in\Bbb R} 


|g_n(y)-g(y)|^{q'}\mes\Omega+\frac{1}{2\varepsilon_n}\bigg| 


\int_\Omega\sup_{y\in\Bbb R}|k_n(x,y)-k(x,y)|dx\bigg|^2. 


\end{gather*} 


Учитывая условие (6), после перехода к пределу приходим к неравенствам 


$$ 


\lim_{n\to\infty}I_n(w_n)\le I(w)\le\inf I(W_\partial)+\delta, 


$$ 


откуда в силу произвольности $\delta$ следует 


\begin{equation} 


\lim_{n\to\infty}I_n(w_n)\le \inf I(W_\partial). 


\end{equation} 


Таким образом, числовая последовательность $\{I_n(w_n)\}$ ограничена. 


Тогда последовательность пар $\{w_n\}$ ограничена в пространстве $S$, 


и справедливы равенства 


\begin{equation} 


\Delta y_n+g_n(y_n)=v_n+g+\varphi_n,\quad \int_\Omega 


k_n(x,y_n)dx=\psi_n, 


\end{equation} 


причем выполняются оценки 


$$ 


\|\varphi_n\|_{q'}\le c(\varepsilon_n)^{1/q'},\quad 


|\psi_n|\le x(\varepsilon_n)^{1/2}. 


$$ 


Учитывая определение множества $G$, установим ограниченность 


последовательности $\{g_n(y_n)\}$ в пространстве $L_{q'}$. Выделяя 


подпоследовательности, получаем сходимости $w_n\to w'=(v',y')$ слабо в 


$S$, $g_n(y_n)\to g'$ 


слабо в $L_{q'}$ и $\varphi_n\to0$ слабо в $L_q$, $\psi_n\to0$, причем 


пара $w'$ является 


элементом множества $S_U$. Пользуясь теоремой 


Реллиха--Кондрашова, после извлечения подпоследовательности 


установим, что $y_n\to y'$ сильно в $L_2$ и п.\,в. на $\Omega$, а значит, 


$g(y_n)\to g(y')$ 


п.\,в. на $\Omega$. Почти для всех $x\in\Omega$ справедливы соотношения 


\begin{gather*} 


|g_n(y_n(x))-g(y'(x))|\le|g_n(y_n(x))-g(y_n(x))|+ 


|g(y_n(x))-g(y'(x))|\le \\ 


% 


\le \sup_{y\in\Bbb R}|g_n(y)-g(y)|+|g(y_n(x))-g(y'(x))|. 


\end{gather*} 


Учитывая условие (3), заключаем, что $g_n(y_n)\to g(y')$ п.\,в. на 


$\Omega$. Применяя 


лемму 1.3 ([2], с.\,25), установим, что $g_n(y_n)\to g(y')$ слабо в 


$L_{q'}$. Умножая 


первое равенство (8) на произвольную достаточно гладкую 


функцию, интегрируя результат по области $\Omega$ и переходя к 


пределу, установим соотношение $\Delta y'+g(y')=v'+f$. Таким образом, 


точка $w'$ 


(слабый предел последовательности $\{w_n\}$ в пространстве $S$), 


оказывается допустимой парой для уравнения (1). 


 


Аналогично почти для всех $x\in\Omega$ устанавливаются неравенства 


\begin{gather*} 


|k_n(x,y_n(x))-k(x,y'(x))|\le|k_n(x,y_n(x))-k(x,y_n(x))|+ 


|k(x,y_n(x))-k(x,y'(x))|\le \\ 


% 


\le\sup_{y\in\Bbb R}|k_n(x,y)-k(x,y)|+ 


|k(x,y_n(x))-k(x,y'(x))|. 


\end{gather*} 


Пользуясь условием (4), заключаем, что $k_n(x,y_n(x))\to k(x,y'(x))$ 


п.\,в. на $\Omega$. 


Применяя теорему Лебега, из второго условия (8) получим 


$$ 


K(y')=\int_\Omega k(x,y')dx=\lim_{n\to\infty}\int_\Omega 


k_n(x,y_n)dx=\lim_{n\to\infty}\psi_n=0. 


$$ 


Отсюда следует, что точка $w'$ принадлежит множеству $W_\partial$. 


 


Из определения функционала $I_n$ и соотношения (7) следует условие 


$$ 


\lim_{n\to\infty}I(w_n)\le\lim_{n\to\infty}I_n(w_n)\le 


\inf I(W_\partial). 


$$ 


Таким образом, для любого сколь угодно малого положительного 


числа $\varepsilon$ номер $n$ можно выбрать столь большим, чтобы значение 


функционала $I(w_n)$ превосходило его нижнюю грань на множестве 


$W_\partial$ не более чем на $\varepsilon$. 


\end{pf} 


 


 


На основе полученных результатов представим алгоритм 


нахождения приближенного (в указанном смысле) решения исходной 


оптимизационной задачи. Для этого достаточно решить в 


соответствии с классическими методами минимизации функционалов 


регуляризованную (разрешимую) задачу. Согласно теореме 3 при 


достаточно больших $n$ величина $w_n$ оказывается слабым 


приближенным решением задачи $P_n$. Характерно, что сходимость 


метода аппроксимации устанавливается даже в отсутствии 


выпуклости функции $F$ и ограничений на степень роста функции $k$, 


применяемых в теореме 1. Тем самым желаемый результат 


получается даже в случае отсутствия оптимального управления в 


задаче $P$. 


 


 


\section{Решение аппроксимационной задачи} 


 


Для вычисления производной минимизируемого функционала в 


задаче $P_n$ потребуются дополнительные ограничения на функции $g_n$ 


и $k_n$. Будем полагать, что они удовлетворяют неравенствам 


$$ 


|g_{ny}(y)|\le a_1+c|y|^{q-2},\quad |k_{ny}(x,y)|\le a_2(x)+ 


c|y|^{q-1}\ \ \forall y\in\Bbb R,\ \ x\in\Omega, 


$$ 


где $a_1>0$, $a_2\in L_{q/(q-1)}(\Omega)$; $g_{ny}(y)$, $k_{ny}(x,y)$ --- 


производные от $g_n$ и $k_n(x,\cdot)$ в точке $y$. 


Тогда операторы $g_n(\cdot):L_q\to L_{q'}$ и $k_n(\cdot):L_q\to L_1$, где 


$k_n(x)=k_n(x,y(x))$, оказываются дифференцируемыми 


по Фреше, причем справедливы равенства ([1], с.\,312) 


$$ 


g'_n(y)h(x)=g_{ny}(y(x))h(x),\quad 


k'_n(y)h(x)=k_{ny}(x,y(x))h(x)\ \ \forall y\in L_q. 


$$ 


Пользуясь теоремой вложения Соболева, установим 


дифференцируемость по Фреше отображения $g_n(\cdot):H^1_0\to H^{-1}$. 


Пусть почти для 


всех $x\in\Omega$ функция $F(x;\cdot,\cdot)$ обладает непрерывными 


производными, причем 


для любых значений $y\in\Bbb R$, $z\in\Bbb R^n$ и п.\,в. на $\Omega$ 


выполняются следующие 


неравенства: $|F(x;y,z)|\le \varphi(x)+ c(|y|^q+|z|^2)$, $|F_y(x;y,z)|\le 


\varphi'(x)+ c(|y|^{q/q'}+ |z|^{2/q'})$, $|F_z(x;y,z)| \le\varphi''(x)+ 


c(|y|^{q/2}+|z|)$, где $\varphi\in L_1$, $\varphi'\in L_{q'}$, 


$\varphi''\in L_2$. Тогда функционал 


$$ 


J(y)=\int_\Omega F(x;y(x),\nabla y(x))dx 


$$ 


будет дифференцируем по Фреше в пространстве $H^1_0$ ([1], с.\,312), 


причем его производная характеризуется равенством 


$$ 


J'(y)h=\int_\Omega[F_y(x;y,\nabla y)h+F_{\nabla y} 


(x;y,\nabla y)\nabla h]dx\ \ \forall h\in H^1_0. 


$$ 


Теперь необходимые условия экстремума для регуляризованной задачи могут 


быть установлены стандартными методами. 


 


\begin{thm} 


Решение $(v_n,y_n)$ задачи $P_n$ характеризуется формулой 


\begin{equation} 


v_n=\Pi(p_n/\gamma) 


\end{equation} 


и уравнениями 


\begin{gather} 


\Delta y_n+g_n(y_n)-v_n-f=(\varepsilon_n)^{q-1} 


|p_n|^{q-2}p_n,\\ 


% 


\Delta p_n+g_{ny}(y_n)p_n=F_y(x;y_n,\nabla y_n)- 


\Div F_{\nabla y} (x;y_n,\nabla y_n)+\varepsilon_n 


\Delta^2y_n+q_n k_{ny}(x,y_n) 


\end{gather} 


с однородными граничными условиями, где $\Pi$ --- проектор на множество 


$U$, причем имеет место равенство 


\begin{equation} 


\int_\Omega k_n(x,y_n(x))dx=\varepsilon_n q_n. 


\end{equation} 


\end{thm} 


 


\begin{pf} 


Функционал $I_n$ в точке $w_n=(v_n,y_n)$ обладает частными производными 


$I_{nv}(w_n)$ и $I_{ny}(w_n)$, определяемыми равенствами 


\begin{gather*} 


I_{nv}(w_n)u=\int_\Omega(\gamma v_n-p_n)u\,dx 


\ \ \forall u\in L_2,\\ 


% 


I_{ny}(w_n)h=\int_\Omega\{F^n_y h+F^n_{\nabla y}\nabla h+ 


\varepsilon_n\Delta y_n\Delta h-[\Delta h+g_{ny}(y_n)h]p_n+ 


q_n k_{ny}(x,y_n)h\}dx\ \ \forall h\in H^1_0, 


\end{gather*} 


где через $F^n_y$ и $F^n_{\nabla y}$ обозначены производные 


$F_y(x;y_n,\nabla y_n)$ и $F_{\nabla y}(x;y_n,\nabla y_n)$, 


\begin{gather*} 


p_n=\frac{1}{\varepsilon_n}|\Delta y_n+g_n(y_n)-v_n-f|^{q'-2} 


[\Delta y_n+g_n(y_n)-v_n-f],\\ 


q_n=\frac{1}{\varepsilon_n}\int_\Omega k_n(x,y_n(x))dx. 


\end{gather*} 


Из последних двух равенств непосредственно следуют соотношения 


(10) и (12). Итак, необходимым условием минимума функционала $I_n$ 


на множестве $S_U$ в точке $w_n$ будут ([18], с.\,18) вариационное 


неравенство 


$$ 


\int_\Omega(\gamma v_n-p_n)(u-v_n)dx\ge0\ \ \forall u\in U 


$$ 


и условие 


$$ 


\int_\Omega\{F^n_y h+F^n_{\nabla y}\nabla h+\varepsilon_n 


\Delta y_n\Delta h-[\Delta h+g_{ny}(y_n)h]p_n+q_n k_{ny} 


(x,y_n)h\}dx=0\ \ \forall h\in H^1_0. 


$$ 


Пользуясь определением оператора проектирования на выпуклое 


подмножество гильбертова пространства, преобразуем 


вариационное неравенство к виду (9). Из последнего соотношения 


в силу формул Грина следует, что функция $p_n$ является решением 


однородной задачи Дирихле для уравнения (11). 


\end{pf} 


 


Согласно теореме 4 решение аппроксимационной задачи сводится к 


рассмотрению однородной задачи Дирихле для системы нелинейных 


эллиптических уравнений 


\begin{gather*} 


\Delta y_n+g_n(y_n)=\Pi(p_n/\gamma)+f+(\varepsilon_n)^{q-1} 


|p_n|^{q-2}p_n,\\ 


\Delta p_n+g_{ny}(y_n)p_n=F_y(x;y_n,\nabla y_n)-\Div 


F_{\nabla y}(x;y_n,\nabla y_n)+\varepsilon_n\Delta^2y_n+ 


q_n k_{ny}(x,y_n) 


\end{gather*} 


и соотношения (12), связывающей неизвестные функции $y_n$, $p_n$ с 


параметром $q_n$. Эта задача может быть решена итерационно. После 


этого определяется функция $v_n$ по формуле (9). В соответствии с 


теоремой 3 при достаточно больших номерах $n$ пара $(v_n,y_n)$ с 


желаемой степенью точности будет удовлетворять уравнению (1), 


а значение функционала $I$ на ней оказывается сколь угодно 


близким к его минимуму на множестве $W_\partial$. Таким образом, 


получаем приближенное решение задачи $P$ в определенном ранее 


смысле. 


 


 


\section{Замечания} 


 


1. Введенное понятие слабого приближенного решения 


оптимизационной задачи аналогично тому, что было определено в 


[13] в отсутствии изопериметрического условия. 


 


2. Аппроксимация негладкого оператора семейством гладких 


осуществлялось в [14] в случае однозначной разрешимости 


уравнения и без изопериметрического условия. 


 


3. Использованный в [3] метод штрафа является исключительно 


средством построения необходимых условий оптимальности для 


исходной задачи. В данном же случае он применяется 


непосредственно для нахождения приближенного решения данной 


задачи, а условия оптимальности для нее вообще не получаются. 


Отметим, что при отсутствии разрешимости задачи эти условия 


вообще не имеют смысла. Отличием используемого выше метода 


штрафа от его известных аналогов является также аппроксимация 


в ``оштрафованных'' членах регуляризованного функционала. 


 


4. Последнее слагаемое в регуляризованном функционале 


напоминает стабилизатор в методе Тихонова (степень нормы с 


малым параметром). Однако обычно регуляризация функционала 


осуществляется с целью преодоления трудностей, обусловленных 


некорректностью задачи в смысле Тихонова [17]. В данном случае 


этот прием позволяет получить дополнительную априорную оценку 


для состояния системы и доказать разрешимость 


аппроксимационной задачи без ограничений, используемых в 


теореме 1. 


 


5. Наличие малого параметра в трех слагаемых регуляризованного 


функционала позволяет получить аппроксимацию исходной 


сингулярной задачи оптимального управления (отсутствие теоремы 


существования и единственности для уравнения состояния, 


негладкость оператора состояния, возможная неразрешимость 


экстремальной задачи) регулярной вариационной задачей 


(разрешимость и гладкость гарантированы). 


 


6. В основе принципа расширения экстремальных задач, 


восходящего к Д.\,Гильберту, лежит переход к расширенной 


задаче так, что минимум функционала для нее совпадает с нижней 


гранью исходного функционала [15], [16]. Хотя полученный 


результат примыкает к теории расширения (ищется нижняя грань 


функционала, которая может оказаться не достижимой), метод 


исследования здесь иной: нижняя грань функционала ищется 


непосредственно на основе аппроксимации задачи $P$ разрешимой 


аппроксимационной задачей. 


 


7. Уравнение (10) представляет собой приближенную форму 


исходного уравнения (1). Действительно, при малых значениях 


параметра $\varepsilon_n$ правая часть этого уравнения будет сколь угодно 


близка к нулю. Аналогично соотношение (12) является 


приближенной формой изопериметрического условия. Это 


согласуется с определением приближенного решения задачи, где 


допускается выполнение ограничений на систему не точно, а 


приближенно. 


 


8. Если сам минимизируемый функционал оказывается не гладким, 


то при определении регуляризованного функционала можно 


провести его гладкую аппроксимацию. Другой вариант 


исследования состоит в использовании описанной выше методики 


совместно с известными методами негладкой оптимизации типа 


[5], [8]. Тогда условие оптимальности для аппроксимационной 


задачи записывается не в терминах обычных функциональных 


производных, а с использованием их обобщений: 


субдифференциала, производной Кларка и т.\,д. 


 


9. При наличии более сложных фазовых ограничений описанную 


методику можно использовать в сочетании с известными 


абстрактными схемами общей теории экстремума [10]--[12]. 


 


10. Степень роста нелинейности в уравнении и определение 


множества $G$ не обязательно связаны с показателем 


интегрируемости $q$, определяемым теоремой Соболева. 


 


 


\begin{thebibliography}{99} 


 


\bibitem{1} 


Крейн С.Г. {\em Функциональный анализ}. -- М.: Наука, 1972. -- 544 с. 


\bibitem{2} 


Лионс Ж.-Л. {\em Некоторые методы решения нелинейных краевых задач}. 


-- М.: Мир, 1972. -- 588 с. 


\bibitem{3} 


Лионс Ж.-Л. {\em Управление сингулярными распределенными системами}. -- 


М.: Наука, 1987. -- 368 с. 


\bibitem{4} 


Фурсиков А.В. {\em Оптимальное управление распределенными системами. 


Теория и приложения}. -- Новосибирск, 1999. 


\bibitem{5} 


Экланд И., Темам Р. {\em Выпуклый анализ и вариационные проблемы}. 


-- М.: Мир, 1979. -- 399 с. 


\bibitem{6} 


Иваненко В.И., Мельник В.С. {\em Вариационные методы в 


задачах управления для систем с распределенными параметрами}. -- 


Киев: Наук. думка, 1988. -- 284 с. 


\bibitem{7} 


Райтум У.Е. {\em Задачи оптимального управления для 


эллиптических уравнений}. -- Рига: Зинатне, 1989. -- 280 с. 


\bibitem{8} 


Серовайский С.Я. {\em Градиентные методы в задаче оптимального управления 


нелинейной эллиптической системой} // Сиб. матем. журн. -- 


1996. -- Т.\,37. -- \No\,5. -- С.\,1154--1166. 


\bibitem{9} 


Кларк Ф. {\em Оптимизация и негладкий анализ}. 


-- М.: Наука, 1988. -- 279 с. 


\bibitem{10} 


Дубовицкий А.О., Милютин А.А. {\em Необходимые условия слабого экстремума 


в общей задаче оптимального управления}. -- 


М.: Наука, 1971. -- 113 с. 


\bibitem{11} 


Иоффе А.Д., Тихомиров В.М. {\em Теория экстремальных задач}. -- М.: 


Наука, 1974. -- 479 с. 


\bibitem{12} 


Якубович В.А. {\em К абстрактной теории оптимального управления} // Сиб. 


матем. журн. -- 1977. -- Т.\,18. -- \No\,3. -- С.\,685--707. 


\bibitem{13} 


Серовайский С.Я. {\em Приближенное решение оптимизационных задач для 


сингулярных бесконечномерных систем} // Сиб. матем. журн. -- 


2003. -- Т.\,44. -- \No\,3. -- С.\,660--673. 


\bibitem{14} 


Серовайский С.Я. {\em Оптимальное управление для уравнений эллиптического 


типа с негладкой нелинейностью} // Дифференц. уравнения. -- 2003. -- 


Т.39. -- \No\,10. -- С.\,1420--1424. 


\bibitem{15} 


Янг Л. {\em Лекции по вариационному исчислению и теории оптимального 


управления}. -- М.: Мир, 1974. -- 488 с. 


\bibitem{16} 


Варга Дж. {\em Оптимальное управление дифференциальными и функциональными 


уравнениями}. -- М.: Мир, 1977. -- 622 с. 


\bibitem{17} 


Васильев Ф.П. {\em Методы решения экстремальных задач. 


Задачи минимизации в функциональных пространствах, регуляризация, 


аппроксимация}. -- М.: Наука, 1981. -- 400 с. 


\bibitem{18} 


Лионс Ж.-Л. {\em Оптимальное управление системами, описываемыми 


уравнениями с частными производными}. -- М.: Мир, 1972. -- 416 


с. 


 


\end{thebibliography} 


 


\vskip 2cm 


 


\post{\it Казахский государственный}{\it Поступила} 


\post{\it национальный университет}{03.06.2003} 


 


\label{end} 


\end{document} 


